\documentclass[11pt,a4paper]{article}

\usepackage[utf8]{inputenc}
\usepackage[T1]{fontenc}
\usepackage[french]{babel}
\usepackage[margin=2.5cm]{geometry}
\usepackage[protrusion=true,expansion=false]{microtype}
\usepackage{booktabs}
\usepackage{xurl}
\usepackage{hyperref}
\usepackage{enumitem}
\usepackage{titlesec}
\usepackage{xcolor}
\usepackage{parskip}
\usepackage{tcolorbox}
\usepackage{tikz}
\usetikzlibrary{arrows.meta, positioning, fit, backgrounds}
\usepackage{fancyhdr}
\usepackage{longtable}
\usepackage{array}
\usepackage{seqsplit}

\tcbuselibrary{skins, breakable}

% --- Palette ---
\definecolor{winkblue}{RGB}{37,60,110}
\definecolor{winkblue2}{RGB}{74,105,170}
\definecolor{codebg}{RGB}{244,245,248}
\definecolor{codeframe}{RGB}{210,214,222}
\definecolor{toolgreen}{RGB}{40,110,80}
\definecolor{warnred}{RGB}{150,40,40}

% --- Code inline : seqsplit permet la coupure a n'importe quel caractere,
% indispensable pour les chemins de fichiers longs sans espace (sinon
% "overfull hbox" garanti dans les tableaux/tikz a largeur fixe). ---
\newcommand{\code}[1]{{\small\ttfamily\color{winkblue}\seqsplit{#1}}}

% --- Titres de section ---
\titleformat{\section}
  {\normalfont\Large\bfseries\color{winkblue}}
  {\thesection}{1em}{}
\titleformat{\subsection}
  {\normalfont\large\bfseries\color{winkblue2}}
  {\thesubsection}{1em}{}
\titlespacing*{\section}{0pt}{1.6em}{0.8em}

% --- En-tête / pied de page ---
\pagestyle{fancy}
\fancyhf{}
\renewcommand{\headrulewidth}{0.4pt}
\fancyhead[L]{\small\textcolor{winkblue}{Wink --- Agent conformité}}
\fancyhead[R]{\small\textcolor{winkblue}{\texttt{compliance\_agent}}}
\fancyfoot[C]{\small\thepage}

\hypersetup{
    colorlinks=true,
    linkcolor=winkblue2,
    urlcolor=winkblue2,
    citecolor=winkblue2,
    pdftitle={Wink -- Agent conformite (compliance\_agent) -- recapitulatif},
    pdfauthor={Wink}
}

% --- Boîte "tool" ---
\newtcolorbox{toolbox}[2]{
  enhanced, breakable,
  colback=white, colframe=toolgreen,
  colbacktitle=toolgreen, coltitle=white,
  fonttitle=\bfseries,
  title={#1},
  arc=2pt, boxrule=0.8pt, left=6pt, right=6pt, top=4pt, bottom=4pt,
  before skip=10pt, after skip=10pt,
  #2
}

% --- Boîte "règle stricte" ---
\newtcolorbox{rulebox}[1][]{
  enhanced, breakable,
  colback=red!4!white, colframe=warnred,
  arc=2pt, boxrule=0.8pt, left=6pt, right=6pt, top=4pt, bottom=4pt,
  before skip=10pt, after skip=10pt,
  title={Garde-fous d'exécution}, fonttitle=\bfseries, coltitle=white,
  colbacktitle=warnred,
  #1
}

% --- Boîte "état actuel / à construire" ---
\newtcolorbox{statebox}[2]{
  enhanced, breakable,
  colback=blue!3!white, colframe=winkblue2,
  arc=2pt, boxrule=0.8pt, left=6pt, right=6pt, top=4pt, bottom=4pt,
  before skip=8pt, after skip=8pt,
  title={#1}, fonttitle=\bfseries, coltitle=white, colbacktitle=winkblue2,
  #2
}

\title{\textcolor{winkblue}{\Huge Wink --- Agent conformité}\\[0.4em]
\Large\textcolor{winkblue2}{\texttt{compliance\_agent} --- Récapitulatif technique}}
\author{}
\date{\today}

\begin{document}

\begin{titlepage}
\centering
\vspace*{3cm}
{\Huge\bfseries\textcolor{winkblue}{Wink}}\\[0.5cm]
{\LARGE\textcolor{winkblue2}{Agent conformité --- \texttt{compliance\_agent}}}\\[0.3cm]
{\large\textcolor{gray}{EPIC-03 --- Conformité \& Normes réglementaires}}\\[2cm]

\begin{tcolorbox}[enhanced, colback=blue!3!white, colframe=winkblue2, arc=3pt,
  boxrule=1pt, width=0.8\textwidth, left=12pt, right=12pt, top=10pt, bottom=10pt]
\centering
\textbf{Objet du document}\\[0.3em]
\small
Rôle actuel de l'agent, ses 3 tools, son mode de raisonnement,
et son point d'intégration prévu avec le futur copilot.\\
À donner tel quel à une nouvelle session pour reprendre ce chantier
sans perdre le contexte.
\end{tcolorbox}

\vfill
{\large\today}
\end{titlepage}

\tableofcontents
\newpage

\section{Contexte et périmètre}

Wink prévoit 3 agents à terme, un par grand domaine métier (décision actée en conversation produit, pas encore reportée dans \code{Wink\_architecture\_fichiers.md}) :

\begin{itemize}[leftmargin=1.2em]
    \item \textbf{\code{administrative\_agent}} (EPIC-01, I1) --- pas encore créé, pas de code.
    \item \textbf{\code{compliance\_agent}} (EPIC-03) --- celui documenté ici, renommé depuis \code{regulatory\_agent}.
    \item \textbf{\code{contract\_agent}} (EPIC-04) --- dossier réservé (\code{agents/contract\_agent/}), vide.
\end{itemize}

\code{compliance\_agent} couvre aujourd'hui 5 tâches~:

\begin{center}
\begin{tabular}{ll}
\toprule
\textbf{Service} & \textbf{Tâche} \\
\midrule
3.3.5 (étape 0) & Détection de secteur \\
3.3.1 & Diagnostic RGPD \\
3.3.12 & Classification AI Act \\
3.3.6 & Diagnostic d'écart HealthTech \\
3.3.6 & Diagnostic d'écart FinTech \\
\bottomrule
\end{tabular}
\end{center}

\begin{tcolorbox}[enhanced, colback=yellow!8!white, colframe=orange!70!black,
  arc=2pt, boxrule=0.8pt, left=6pt, right=6pt, top=4pt, bottom=4pt]
\textbf{Seule la tâche HealthTech gap-analysis a été convertie} en agent à raisonnement autonome (tool-calling) --- c'est le \textbf{pilote}. RGPD, AI Act, FinTech gap-analysis et détection de secteur restent sur l'ancien pattern (un seul appel LLM, RAG pré-effectué par le code appelant, pas de tools). Conversion prévue service par service, jamais en une fois.
\end{tcolorbox}

\section{Ce que fait l'agent (pilote HealthTech, service 3.3.6)}

\subsection{Avant / Maintenant}

\begin{center}
\begin{tikzpicture}[
  node distance=0.9cm and 1.4cm,
  box/.style={rectangle, rounded corners, draw=winkblue2, fill=blue!4!white,
              text width=4.6cm, align=center, minimum height=1.1cm, font=\small},
  lbl/.style={font=\bfseries\small\color{winkblue}}
]
  \node[lbl] (avantlbl) {Avant (ancien pattern)};
  \node[box, below=of avantlbl] (a1) {Code Python calcule l'embedding};
  \node[box, below=of a1] (a2) {Code Python cherche les chunks};
  \node[box, below=of a2] (a3) {Code Python récupère la classification};
  \node[box, below=of a3, fill=orange!10, draw=orange!70!black] (a4) {LLM : un seul appel, rédige le JSON};

  \node[lbl, right=2.6cm of avantlbl] (maintlbl) {Maintenant (pilote HealthTech)};
  \node[box, below=of maintlbl, fill=green!8, draw=toolgreen] (m1) {Agent reçoit la tâche};
  \node[box, below=of m1, fill=green!8, draw=toolgreen] (m2) {Agent décide : appeler un tool ?};
  \node[box, below=of m2, fill=green!8, draw=toolgreen] (m3) {Agent recherche / relit si besoin};
  \node[box, below=of m3, fill=green!8, draw=toolgreen] (m4) {Agent produit la réponse finale};

  \draw[-{Latex[length=2mm]}, winkblue2] (a1) -- (a2);
  \draw[-{Latex[length=2mm]}, winkblue2] (a2) -- (a3);
  \draw[-{Latex[length=2mm]}, winkblue2] (a3) -- (a4);

  \draw[-{Latex[length=2mm]}, toolgreen] (m1) -- (m2);
  \draw[-{Latex[length=2mm]}, toolgreen] (m2) -- (m3);
  \draw[-{Latex[length=2mm]}, toolgreen] (m3) -- (m4);
  \draw[-{Latex[length=2mm]}, toolgreen, dashed] (m3.east) to[out=0,in=0,looseness=2.2] (m2.east);
\end{tikzpicture}
\end{center}

Le code Python ne décide plus rien à l'avance~: il donne à l'agent le chapitre à traiter et une requête de recherche suggérée, puis le laisse \textbf{décider lui-même} s'il doit consulter le profil du projet, rechercher dans le corpus légal, combien de fois, avec quelle formulation.

\subsection{Un run typique}

\begin{enumerate}[leftmargin=1.4em]
    \item L'agent reçoit la tâche (chapitre à traiter + requête suggérée).
    \item Il appelle \code{lire\_profil\_projet} pour connaître la classe de risque MDR du projet.
    \item Il appelle \code{rechercher\_corpus\_reglementaire} avec la requête suggérée (ou une reformulation).
    \item Si les résultats lui semblent insuffisants, il peut rechercher à nouveau avec une autre requête.
    \item Il produit la liste d'obligations, chacune citant le \code{source\_id} réel d'un chunk qu'il a vu.
\end{enumerate}

\section{Les 3 tools}
\label{sec:tools}

Tous en \textbf{lecture seule} --- l'agent ne persiste jamais rien lui-même, ça reste le rôle du service backend appelant (\code{upsert\_result}). Définis dans \code{backend/entrepreneur/services/compliance\_agent\_tools\_service.py} (pas dans \code{agents/compliance\_agent/}, voir section~\ref{sec:separation}).

\begin{toolbox}{1. \texttt{rechercher\_corpus\_reglementaire(domaine, requete)}}{}
Recherche sémantique dans le corpus légal interne déjà ingéré (\code{content\_embeddings}), filtrée par domaine (\code{secteur\_healthtech}, \code{secteur\_fintech}, \code{rgpd}, \code{ai\_act}\ldots).

\textbf{Renvoie} : une liste de chunks \texttt{\{source\_id, texte\}}.

\textbf{En interne} : \code{embed\_texts} (embedding de la requête) puis \code{search\_similar\_chunks} (recherche cosinus, top~8).
\end{toolbox}

\begin{toolbox}{2. \texttt{lire\_profil\_projet()}}{}
Contexte réglementaire ET business déjà connu du projet : secteur détecté, classification (classe de risque MDR ou type d'agrément FinTech) si elle existe, réponses au questionnaire sectoriel, et champs business du projet (stage, description, technologies, customer\_segments, target\_users, value\_proposition).

\textbf{Sert à} : contextualiser une génération.
\end{toolbox}

\begin{toolbox}{3. \texttt{lire\_diagnostics\_existants()}}{}
Résultats déjà persistés pour ce projet (obligations RGPD, classification AI Act, chapitres de gap-analysis HealthTech/FinTech avec leurs obligations et le statut manuel « accompli »).

\textbf{Sert à} : répondre sans régénérer --- pensé pour le copilot (section~\ref{sec:copilot}), pas encore utilisé par le pilote actuel qui ne fait que générer.
\end{toolbox}

\begin{tcolorbox}[enhanced, colback=red!4!white, colframe=warnred, arc=2pt,
  boxrule=0.8pt, left=6pt, right=6pt, top=4pt, bottom=4pt,
  title={Explicitement exclus}, fonttitle=\bfseries, coltitle=white, colbacktitle=warnred]
\begin{itemize}[leftmargin=1.2em]
  \item Un \textbf{tool d'écriture} --- l'agent ne modifie jamais la DB (frontière de sécurité, pas une limite technique).
  \item Un \textbf{tool de recherche web} --- casserait la garantie de source vérifiable (voir section~\ref{sec:raisonnement}).
\end{itemize}
\end{tcolorbox}

\section{Comment il raisonne, et garde-fous}
\label{sec:raisonnement}

Implémenté avec \textbf{pydantic-ai} (\code{agents/compliance\_agent/healthtech\_gap\_analysis\_service.py}) : un \code{pydantic\_ai.Agent} construit avec un modèle Groq (\code{openai/gpt-oss-120b}), un system prompt, les 3~tools, et un schéma de sortie structuré (\code{GapAnalysisResult}).

\subsection{Règles imposées par le system prompt}

\begin{itemize}[leftmargin=1.2em]
    \item Jamais d'obligation inventée sans chunk réellement vu pendant la tâche en cours.
    \item Les obligations doivent tenir compte de la classe de risque trouvée via \code{lire\_profil\_projet}.
    \item Liste vide acceptée et correcte si le corpus ne couvre pas le sujet --- \textbf{jamais de recherche web de secours}. Un trou dans le corpus doit remonter pour un enrichissement contrôlé, pas être patché silencieusement à chaque appel.
    \item Peut suggérer textuellement qu'une obligation semble déjà couverte par un fait du profil, mais ne coche jamais lui-même une case « accompli » --- action manuelle de l'entrepreneur.
    \item \code{source\_id} doit être un identifiant réel vu via le tool de recherche, jamais inventé.
\end{itemize}

\begin{rulebox}
\begin{itemize}[leftmargin=1.2em]
  \item \texttt{UsageLimits(request\_limit=10, tool\_calls\_limit=6)} --- borne le nombre d'appels pour éviter qu'un agent jamais satisfait de ses recherches ne boucle indéfiniment (chaque tour coûte un appel LLM réel).
  \item \textbf{Vérification a posteriori} (\code{\_obligations\_to\_dict}, côté backend) : chaque \code{source\_id} renvoyé par l'agent est comparé à l'ensemble des \code{source\_id} réellement vus pendant ce run. Un \code{source\_id} absent (hallucination) devient \code{None}, pas une erreur bloquante.
\end{itemize}
\end{rulebox}

\subsection{Comment il choisit quel tool appeler}

C'est le LLM lui-même, pas du code. Le system prompt décrit ce que fait chaque tool (via sa signature et sa docstring, exposées automatiquement par pydantic-ai comme schéma de function-calling) ; à chaque étape, le modèle Groq décide s'il a besoin d'un tool, lequel, avec quels arguments, en se basant sur la tâche reçue et les résultats des appels précédents --- le mécanisme de function-calling natif de Groq (API compatible OpenAI).

\section{Pourquoi les tools ne sont pas dans \texorpdfstring{\texttt{agents/}}{agents/}}
\label{sec:separation}

Règle déjà actée avant ce chantier : \textbf{\code{agents/} ne touche jamais la DB directement} (\code{backend/} possède SQLAlchemy et les \code{Session}).

\begin{center}
\begin{tikzpicture}[
  node distance=1.4cm,
  box/.style={rectangle, rounded corners, draw=winkblue2, fill=blue!4!white,
              text width=4.4cm, align=center, minimum height=1.3cm, font=\small}
]
  \node[box, fill=green!8, draw=toolgreen] (backend) {\textbf{backend/}\\possède SQLAlchemy, les \texttt{Session}, la persistance};
  \node[box, right=of backend] (agents) {\textbf{agents/}\\prompts, schémas Pydantic, appels LLM --- jamais de DB};

  \draw[-{Latex[length=3mm]}, winkblue, thick] (backend) -- node[above, font=\small]{importe} (agents);
  \draw[-{Latex[length=3mm]}, warnred, thick, dashed] ([yshift=-0.5cm]agents.south) to[out=-90,in=-90, looseness=1.3]
    node[below, font=\small\color{warnred}]{jamais l'inverse} ([yshift=-0.5cm]backend.south);
\end{tikzpicture}
\end{center}

Deux des trois tools ont besoin d'une \code{Session} pour lire la DB --- donc leur implémentation vit dans \code{backend/entrepreneur/services/compliance\_agent\_tools\_service.py}, fermée par closure sur \texttt{(db, user, project\_id)}, puis \textbf{injectée} à l'agent au moment de l'appel :

\begin{center}
\texttt{generate\_gap\_analysis(chapitre\_titre, requete, tools)}
\end{center}

Cela respecte le sens de dépendance déjà en place sans dépendance circulaire. Le fichier \code{agents/compliance\_agent/healthtech\_gap\_analysis\_service.py} ne sait pas que ces fonctions touchent une DB --- il reçoit juste une liste de fonctions Python à enregistrer sur l'\code{Agent}. L'endroit où vit le code d'un tool ne change rien au fait que c'est bien le LLM qui décide de l'appeler : c'est une question d'organisation de dépôt, pas de définition d'un agent.

\section{Rôle prévu dans le futur copilot --- pas encore construit}
\label{sec:copilot}

\begin{statebox}{État actuel --- ce qui existe réellement}{}
\code{compliance\_agent} n'est appelé QUE depuis la chaîne suivante :
\end{statebox}

\begin{center}
\begin{tikzpicture}[
  node distance=0.55cm,
  box/.style={rectangle, rounded corners, draw=winkblue2, fill=blue!4!white,
              text width=8.6cm, align=left, minimum height=0.95cm, font=\small, inner sep=6pt}
]
  \node[box] (c1) {\textbf{Clic} « Générer »/« Régénérer » --- \code{SectorDiagnosticPage.tsx}};
  \node[box, below=of c1] (c2) {\textbf{Route} \texttt{\seqsplit{POST /projects/\{id\}/sector/healthtech/gap-analysis/\{chapitre\}}} --- \code{sector\_router.py}};
  \node[box, below=of c2] (c3) {\textbf{Service} \code{run\_chapter\_diagnostic} --- \code{healthtech\_gap\_analysis\_diagnostic\_service.py}};
  \node[box, below=of c3, fill=green!8, draw=toolgreen] (c4) {\textbf{Agent} \code{compliance\_agent} (boucle de raisonnement à tools)};

  \draw[-{Latex[length=2.5mm]}, winkblue, thick] (c1) -- (c2);
  \draw[-{Latex[length=2.5mm]}, winkblue, thick] (c2) -- (c3);
  \draw[-{Latex[length=2.5mm]}, winkblue, thick] (c3) -- (c4);
\end{tikzpicture}
\end{center}

Aucun copilot n'existe encore --- \code{agents/orchestrator/} est un dossier vide (juste un \code{.gitkeep}), aucune route de conversation libre n'est câblée.

\begin{statebox}{Ce qui est prévu --- décision produit, pas encore implémenté}{}
Le futur copilot entrepreneur posera des questions libres (« est-ce que je dois faire une AIPD~? », « où en est mon dossier MDR~? »).
\end{statebox}

\begin{center}
\begin{tikzpicture}[
  node distance=0.55cm,
  box/.style={rectangle, rounded corners, draw=gray, fill=gray!6,
              text width=8.6cm, align=left, minimum height=0.95cm, font=\small, inner sep=6pt, dashed}
]
  \node[box] (f1) {\textbf{Question libre} de l'entrepreneur (copilot --- pas encore construit)};
  \node[box, below=of f1] (f2) {\textbf{Orchestrateur} \code{agents/orchestrator/} (à construire) --- route vers le bon agent};
  \node[box, below=of f2, draw=toolgreen] (f3) {\textbf{Agent} \code{compliance\_agent} --- utilise \code{lire\_diagnostics\_existants} pour répondre sans régénérer};

  \draw[-{Latex[length=2.5mm]}, gray, thick, dashed] (f1) -- (f2);
  \draw[-{Latex[length=2.5mm]}, gray, thick, dashed] (f2) -- (f3);
\end{tikzpicture}
\end{center}

C'est précisément pour cet usage que \code{lire\_diagnostics\_existants} a été conçu dès maintenant, même si le pilote actuel (génération pure) ne s'en sert pas encore. Reste à construire : la route de conversation libre, l'orchestrateur qui décide vers quel agent router, et la conversion des 4 autres tâches (RGPD, AI Act, FinTech gap, détection secteur) au même pattern à tools que le pilote HealthTech.

\section{Fichiers touchés}

\begin{longtable}{>{\raggedright\arraybackslash}p{6.7cm}>{\raggedright\arraybackslash}p{8cm}}
\toprule
\textbf{Fichier} & \textbf{Changement} \\
\midrule
\endhead
\code{agents/compliance\_agent/} & Renommé depuis \code{regulatory\_agent/} (5 fichiers : \texttt{ai\_act\_service.py}, \texttt{fintech\_gap\_analysis\_service.py}, \texttt{healthtech\_gap\_analysis\_service.py} [converti en agent], \texttt{rgpd\_service.py}, \texttt{sector\_detection\_service.py}) \\[4pt]
\code{compliance\_agent\_tools\_service.py} & Nouveau fichier (\code{backend/entrepreneur/services/}) --- les 3 tools, accès DB \\[4pt]
\code{healthtech\_gap\_analysis\_diagnostic\_service.py} & Réécrit pour construire les tools et appeler l'agent au lieu de faire le RAG lui-même \\[4pt]
\code{llm\_config.yaml} / \code{prompts.yaml} & Clés renommées \texttt{regulatory\_agent\_*} $\rightarrow$ \texttt{compliance\_agent\_*} ; nouvelle clé \texttt{system} (remplace \texttt{diagnostic}) pour \texttt{compliance\_agent\_healthtech\_gap} \\[4pt]
\code{backend/requirements.txt} & Ajout \texttt{pydantic-ai-slim[groq]>=2.46.0}, bump \texttt{pydantic>=2.12.0} et \texttt{groq>=0.25.0} \\[4pt]
\texttt{RegulatoryAgentError} & Renommé \texttt{RgpdAgentError} (cohérence avec les autres exceptions du dossier) \\
\bottomrule
\end{longtable}

\begin{tcolorbox}[enhanced, colback=yellow!8!white, colframe=orange!70!black,
  arc=2pt, boxrule=0.8pt, left=6pt, right=6pt, top=4pt, bottom=4pt,
  title={Non testé de bout en bout}, fonttitle=\bfseries, coltitle=black, colbacktitle=orange!30!white]
\code{GROQ\_API\_KEY} existe déjà dans \code{backend/.env} (ajoutée pour l'extraction de documents, \code{agents/shared/document\_processing/}) et sera réutilisée telle quelle par \code{GroqProvider} --- mais aucun vrai appel Groq avec la boucle de tool-calling n'a encore été exécuté sur ce chantier : seule la syntaxe et la résolution des imports ont été vérifiées (\code{py\_compile}).
\end{tcolorbox}

\section{Marche à suivre pour tester}

\subsection{Préparation}
\begin{enumerate}[leftmargin=1.4em]
  \item \code{GROQ\_API\_KEY} est déjà présente dans \code{backend/.env} --- rien à faire.
  \item Installer les nouvelles dépendances dans le \textbf{vrai} environnement du projet (\code{.venv} Windows) : \texttt{pip install -r backend/requirements.txt}.
\end{enumerate}

\subsection{Test technique --- ça tourne sans erreur ?}
Le plus rapide : lancer le backend (\code{uvicorn}), ouvrir un projet HealthTech au questionnaire complet, cliquer « Générer » sur le premier chapitre, regarder les logs serveur.

Pour un test isolé sans passer par l'UI :
\begin{tcolorbox}[enhanced, breakable, colback=codebg, colframe=codeframe,
  arc=2pt, boxrule=0.6pt, left=6pt, right=6pt, top=4pt, bottom=4pt]
\small\ttfamily
from backend.shared.database import SessionLocal\\
from backend.entrepreneur.services.compliance\_agent\_tools\_service import build\_compliance\_agent\_tools\\
from agents.compliance\_agent.healthtech\_gap\_analysis\_service import generate\_gap\_analysis\\[4pt]
db = SessionLocal()\\
seen = []\\
tools = build\_compliance\_agent\_tools(db, user, project\_id, seen)\\
result = generate\_gap\_analysis("Obligations du fabricant", "obligations du fabricant avant mise sur le marché MDR", tools)\\
print(result.model\_dump\_json(indent=2))\\
print("Chunks vus :", seen)
\end{tcolorbox}

\subsection{Validation de fond --- est-ce correct sur le plan réglementaire ?}

\begin{itemize}[leftmargin=1.2em]
  \item \textbf{Traçabilité} : logue systématiquement les \code{seen\_source\_ids} à côté du résultat --- vérifie que le contenu cité est pertinent, pas juste que l'ID existe.
  \item \textbf{Relecture humaine} : pour chaque \code{source\_id} cité, relire le chunk réel (ou l'article MDR correspondant sur EUR-Lex) et vérifier que l'obligation générée en est une interprétation fidèle. Aucun test automatisé ne remplace ça pour les premiers runs de chaque chapitre.
  \item \textbf{Sensibilité à la classe de risque} : générer le même chapitre pour deux projets de classes différentes (I vs III) --- les obligations doivent différer en conséquence.
  \item \textbf{Cas limites} : classe de risque non déterminée (ne doit rien halluciner), domaine sans chunk pertinent (doit renvoyer une liste vide justifiée), rejouer deux fois le même chapitre à température~0 (résultats globalement stables).
  \item \textbf{Consommation} : \code{result.usage()} donne le nombre de requêtes/tool-calls réellement utilisés --- doit rester loin des limites (6 tool-calls, 10 requêtes) ; sinon, l'agent cherche trop ou le corpus est mal indexé pour ce chapitre.
\end{itemize}

\end{document}
